
\documentclass{article}
\usepackage{colortbl}
\usepackage{makecell}
\usepackage{multirow}
\usepackage{supertabular}

\begin{document}

\newcounter{utterance}

\centering \large Interaction Transcript for game `cladder', experiment `full\_v1.5\_default', episode 1743 with Qwen3.5{-}9B.
\vspace{24pt}

{ \footnotesize  \setcounter{utterance}{1}
\setlength{\tabcolsep}{0pt}
\begin{supertabular}{c@{$\;$}|p{.15\linewidth}@{}p{.15\linewidth}p{.15\linewidth}p{.15\linewidth}p{.15\linewidth}p{.15\linewidth}}
    \# & \multicolumn{2}{c}{Player} && \multicolumn{2}{c}{Game Master} \\
    \hline

    \theutterance \stepcounter{utterance}  
    & & & \multicolumn{4}{p{0.6\linewidth}}{
        \cellcolor[rgb]{0.9,0.9,0.9}{
            \makecell[{{p{\linewidth}}}]{
                \texttt{\tiny{[P1$\langle$GM]}}
                \texttt{You are an expert in causal inference. The following question is not a typical commonsense query, but rather a meticulously designed question created by a professor specializing in causal inference, intended to assess the students' mastery of the course content.} \\
\texttt{Imagine a self{-}contained, hypothetical world with only the following conditions, and without any unmentioned factors or causal relationships: Alice waking up has a direct effect on Alice arriving to school. Traffic has a direct effect on Alice arriving to school. For days when Alice wakes up on time, the probability of arriving to school on time is 46\%. For days when Alice wakes up late, the probability of arriving to school on time is 78\%. For days when Alice wakes up late, would it be more likely to see arriving to school on time if Alice had gotten up on time?} \\
\texttt{Start your answer with 'Yes' or 'No', followed by additional reasoning or evidence to support your explanation.} \\
            }
        }
    }
     \\ \\

    \theutterance \stepcounter{utterance}  
    & \multicolumn{4}{p{0.6\linewidth}}{
        \cellcolor[rgb]{0.561,0.737,0.561}{
            \makecell[{{p{\linewidth}}}]{
                \texttt{\tiny{[P1$\rangle$GM]}}
                \texttt{No} \\
\\ 
\texttt{**Reasoning:**} \\
\\ 
\texttt{To determine the answer, we must analyze the causal structure and the specific probabilities provided within the hypothetical world.} \\
\\ 
\texttt{1.  **Causal Structure**: The problem states that "Alice waking up" has a direct effect on "Alice arriving to school." This establishes a causal link where the state of waking up (On Time vs. Late) is the cause, and the state of arriving (On Time vs. Late) is the effect.} \\
\texttt{2.  **Counterfactual Definition**: The question asks: "For days when Alice wakes up late, would it be more likely to see arriving to school on time if Alice had gotten up on time?"} \\
\texttt{    *   **Observed Reality (\$Y\_\{obs\}\$)**: Given that Alice woke up late (\$W=Late\$), the probability of arriving on time is explicitly stated as **78\%**.} \\
\texttt{    *   **Counterfactual Scenario (\$Y\_\{cf\}\$)**: We are asked to evaluate the probability of arriving on time (\$Y=OnTime\$) under the condition that Alice *had* gotten up on time (\$W=OnTime\$), specifically for the same days where she actually woke up late.} \\
\texttt{3.  **Logical Contradiction in the Hypothetical World**: The prompt defines a world with "only the following conditions" and "without any unmentioned factors."} \\
\texttt{    *   If Alice wakes up late, she is in the state \$W=Late\$.} \\
\texttt{    *   The counterfactual asks what happens if she *had* gotten up on time (\$W=OnTime\$).} \\
\texttt{    *   In a deterministic or strictly defined causal model without unmentioned confounders, the state of "waking up" is a binary event. If the actual event was "waking up late," the counterfactual event "waking up on time" represents a different trajectory entirely.} \\
\texttt{    *   However, the crucial logical trap lies in the phrasing of the probabilities. The prompt gives us \$P(Y=OnTime | W=Late) = 0.78\$. It does *not* give us \$P(Y=OnTime | W=OnTime)\$. It only gives us the probability for the "On Time" waking scenario in the context of the first sentence: "For days when Alice wakes up on time, the probability of arriving to school on time is 46\%."} \\
\texttt{    *   Therefore, \$P(Y=OnTime | W=OnTime) = 0.46\$.} \\
\texttt{4.  **Comparison**:} \\
\texttt{    *   Probability of arriving on time given she woke up late (Actual): **78\%**} \\
\texttt{    *   Probability of arriving on time given she woke up on time (Counterfactual): **46\%**} \\
\texttt{5.  **Conclusion**: The question asks if it would be *more* likely to arrive on time in the counterfactual} \\
            }
        }
    }
    & & \\ \\

    \theutterance \stepcounter{utterance}  
    & & & \multicolumn{2}{p{0.3\linewidth}}{
        \cellcolor[rgb]{0.9,0.9,0.9}{
            \makecell[{{p{\linewidth}}}]{
                \texttt{\tiny{[GM$|$GM]}}
                \texttt{No} \\
            }
        }
    }
    & & \\ \\

    \theutterance \stepcounter{utterance}  
    & & & \multicolumn{2}{p{0.3\linewidth}}{
        \cellcolor[rgb]{0.9,0.9,0.9}{
            \makecell[{{p{\linewidth}}}]{
                \texttt{\tiny{[GM$|$GM]}}
                \texttt{no} \\
            }
        }
    }
    & & \\ \\

    \theutterance \stepcounter{utterance}  
    & & & \multicolumn{2}{p{0.3\linewidth}}{
        \cellcolor[rgb]{0.9,0.9,0.9}{
            \makecell[{{p{\linewidth}}}]{
                \texttt{\tiny{[GM$|$GM]}}
                \texttt{game\_result = WIN} \\
            }
        }
    }
    & & \\ \\

\end{supertabular}
}

\end{document}
